\section{Limitations and Future Work}

Our current study focuses on validating ReSync under the experimental settings available within our compute budget: \textbf{[PLACEHOLDER: completed datasets/modalities]}. Due to the cost of large-scale generative pretraining, we do not yet evaluate ReSync on open-domain text-to-video pretraining or broad multimodal training at the same scale as recent frontier models. The final scope of completed runs, compute budget, and unfinished experiments should be reported explicitly: \textbf{[PLACEHOLDER: final compute budget and unfinished runs]}. Accordingly, any claims about video or large-scale multimodal generation should be viewed as future directions rather than established empirical conclusions in this submission.

We believe the extension to video is conceptually natural because ReSync is self-contained: it uses the diffusion transformer's own internal depth trajectory rather than relying on a strong external representation encoder. This property is particularly relevant for domains where pretrained encoders may be weaker, less standardized, or mismatched to the target generative distribution. However, this hypothesis requires direct validation on video data, including comparisons against external-representation alignment methods and self-contained alignment baselines.

ReSync also suggests a possible path toward layer-skipping inference through output distillation. In this paper, we treat this as an auxiliary capability rather than the main claim. A complete evaluation should measure both generation quality and wall-clock speed under different skip lengths, predictor sizes, and output-distillation weights: \textbf{[PLACEHOLDER: skip-inference ablation]}. Future work should further study when predicted intermediate states remain stable under long skips and whether transition consistency reduces sensitivity to layer-pair selection.

\section{Conclusion}

We introduced ReSync, a self-contained objective for learning compositional internal representation transitions in Diffusion Transformers. Rather than directly aligning two hidden layers, ReSync predicts how an earlier representation should transform into a later one and enforces consistency between composed and direct transitions across depth. This framing combines the simplicity of internal self-alignment with a transition model that can support representation learning and future layer-skipping inference through optional output distillation.
